% Source: https://artofproblemsolving.com/community/c3h1967487_equation_in_nonzerointeger_set
% Stable identity: aops-topic-1967487-q01
Multiplying by $xy\ne0$ gives $2+3y+4x=-5xy$, that is $5xy+4x+3y+2=0$. Multiply by $5$ and factor (Simon's trick):
\[
25xy+20x+15y+10=0\quad\Longleftrightarrow\quad(5x+3)(5y+4)=12-10=2 .
\]
So $5x+3$ is a divisor of $2$, i.e. one of $\pm1,\pm2$, and it is $\equiv3\pmod5$. Only $-2$ qualifies, giving $x=-1$ and $5y+4=-1$, $y=-1$. Check: $\dfrac{2}{1}-3-4=-5$. The only solution is $\boxed{(x,y)=(-1,-1)}$.
